\documentclass[../../main.tex]{subfiles}
% \documentclass[a4paper,10pt]{ltjsarticle}
\usepackage{luatexja}

\usepackage[truedimen,top=25truemm,bottom=20truemm,left=15truemm,right=15truemm]{geometry}
\setlength{\textwidth}{54\zw}
\setlength{\textheight}{73\zw}

\usepackage{amsmath,amssymb,ascmac}
\usepackage{enumerate}
\usepackage{multicol}
\usepackage{cleveref}
\usepackage{framed}
\usepackage{fancyhdr}
\usepackage{latexsym}
\usepackage{tikz}
\usepackage{pgfplots}
\usetikzlibrary{math}
% \usepackage{indent}
\usepackage{cases}
\usepackage{float}
\allowdisplaybreaks
\pagestyle{fancy}
\lhead{}
\chead{}
\begin{document}
\rhead{京都大学後期$1999$年$4$番}
\setcounter{equation}{0}
\begin{oframed}
  \subfile{problem}
\end{oframed}
\setlength{\columnseprule}{0.4pt}
\begin{multicols}{2}
  {\bf [解答]}

  \begin{figure}[H]
    \centering
    \begin{tikzpicture}[scale=0.85, every node/.style={font=\small}]
      % 図1と同じ例 a^2=14.24, b^2=9.76, c^2=16 （a=BC, b=CA, c=AB）
      \pgfmathsetmacro{\ra}{sqrt(14.24)}
      \pgfmathsetmacro{\rb}{sqrt(9.76)}
      \pgfmathsetmacro{\angC}{acos((14.24+9.76-16)/(2*\ra*\rb))}
      \coordinate (C) at (0,0);
      \coordinate (B) at (\ra,0);
      \coordinate (A) at (\angC:\rb);
      % 余弦定理で用いる角 ACB
      \draw[thick, gray!70!black] ($(C)!0.55cm!(B)$) arc[start angle=0, end angle=\angC, radius=0.55cm];
      \draw[line width=1.3pt, red!85!black] (B) -- (C);
      \draw[line width=1.3pt, blue!85!black] (C) -- (A);
      \draw[line width=1.3pt, teal!85!black] (A) -- (B);
      \node[red!85!black, below=1pt] at ($(B)!0.5!(C)$) {$a$};
      \node[blue!85!black, left=1pt] at ($(C)!0.5!(A)$) {$b$};
      \node[teal!85!black, above right=0pt] at ($(A)!0.5!(B)$) {$c$};
      \foreach \p/\pos in {A/above, B/below right, C/below left} {
        \node[circle, fill=black, inner sep=1.2pt] at (\p) {};
        \node[\pos=1pt] at (\p) {$\mathrm{\p}$};
      }
    \end{tikzpicture}
    \caption{鋭角三角形$\mathrm{ABC}$（$a=\mathrm{BC}$，$b=\mathrm{CA}$，$c=\mathrm{AB}$）}
    \label{1999-4:fig:triangle}
  \end{figure}

  $\triangle ABC$の$3$辺の長さを$|BC|=a,|CA|=b,|AB|=c$とする ($a,b,c > 0$)．
  $\triangle ABC$が鋭角三角形だから，余弦定理より
  \begin{align}
    c^2 = a^2+b^2 - 2ab\cos\angle\mathrm{ACB} < a^2+b^2
  \end{align}
  のような式が成り立つから，
  \begin{align}
    \begin{cases}
      x^2 = \dfrac{a^2+b^2-c^2}{2} > 0 \\
      y^2 = \dfrac{a^2-b^2+c^2}{2} > 0 \\
      z^2 = \dfrac{-a^2+b^2+c^2}{2} > 0
    \end{cases}
  \end{align}
  となる$x,y,z \in \mathbb{R}^+$が存在する．
  従って，これら$x,y,z$を$3$辺の長さとする直方体が存在する．
  この時図のように四面体$ABCD$をとると各面が辺の長さ$a,b,c$の三角形で合同になっているから題意は示された．$\cdots$(答)

  \begin{figure}[H]
    \centering
    \begin{tikzpicture}[x={(-0.62cm,-0.5cm)}, y={(1.05cm,0cm)}, z={(0cm,1.05cm)},
      every node/.style={font=\small}]
      % 直方体の辺 x,y,z（図示用の値）．a^2=x^2+y^2, b^2=x^2+z^2, c^2=y^2+z^2
      % x=2, y=3.2, z=2.4 のとき a^2=14.24, b^2=9.76, c^2=16 で鋭角三角形になる
      \def\X{2}
      \def\Y{3.2}
      \def\Z{2.4}
      % 斜投影で手前は x 大・上は z 大・右は y 大．原点 O が隠れた頂点
      \coordinate (O)   at (0,0,0);
      \coordinate (Px)  at (\X,0,0);
      \coordinate (Py)  at (0,\Y,0);
      \coordinate (Pz)  at (0,0,\Z);
      \coordinate (Pxy) at (\X,\Y,0);
      \coordinate (Pxz) at (\X,0,\Z);
      \coordinate (Pyz) at (0,\Y,\Z);
      \coordinate (A)   at (\X,\Y,\Z);
      \coordinate (B)   at (Pz);
      \coordinate (C)   at (Py);
      \coordinate (D)   at (Px);

      % 直方体：隠れた稜線（O に接する3本）は破線
      \draw[densely dashed, gray!70!black] (O) -- (Px);
      \draw[densely dashed, gray!70!black] (O) -- (Py);
      \draw[densely dashed, gray!70!black] (O) -- (Pz);
      \draw[thick, gray!80!black] (Px) -- (Pxy) -- (Py) -- (Pyz) -- (Pz) -- (Pxz) -- cycle;
      \draw[thick, gray!80!black] (Pxz) -- (A) -- (Pyz);
      \draw[thick, gray!80!black] (A) -- (Pxy);

      % 四面体 ABCD：向かい合う辺は同じ長さ（同色）．裏面上の辺は破線
      \draw[line width=1.3pt, red!85!black] (A) -- (B);
      \draw[line width=1.3pt, red!85!black, dashed] (C) -- (D);
      \draw[line width=1.3pt, blue!85!black] (A) -- (C);
      \draw[line width=1.3pt, blue!85!black, dashed] (B) -- (D);
      \draw[line width=1.3pt, teal!85!black] (A) -- (D);
      \draw[line width=1.3pt, teal!85!black, dashed] (B) -- (C);

      % 四面体の辺の長さ
      \node[red!85!black, above=1pt] at ($(A)!0.5!(B)$) {$a$};
      \node[red!85!black, below=1pt] at ($(C)!0.62!(D)$) {$a$};
      \node[blue!85!black, right=1pt] at ($(A)!0.45!(C)$) {$b$};
      \node[blue!85!black, left=1pt] at ($(B)!0.55!(D)$) {$b$};
      \node[teal!85!black, left=1pt] at ($(A)!0.42!(D)$) {$c$};
      \node[teal!85!black, below left=0pt] at ($(B)!0.42!(C)$) {$c$};

      % 直方体の辺の長さ
      \node[below left=0pt] at ($(Pxy)!0.5!(Py)$) {$x$};
      \node[below=1pt] at ($(Px)!0.5!(Pxy)$) {$y$};
      \node[right=1pt] at ($(Pxy)!0.5!(A)$) {$z$};

      % 頂点
      \foreach \p/\pos in {A/above right, B/above left, C/right, D/below left} {
        \node[circle, fill, inner sep=1.3pt] at (\p) {};
        \node[\pos=1pt] at (\p) {$\mathrm{\p}$};
      }
    \end{tikzpicture}
    \caption{題意を満たす直方体と四面体（向かい合う辺の長さが等しい）}
    \label{1999-4:fig:1}
  \end{figure}

  \vspace{10pt}
  {\bf [解説]}

  \newpage
\end{multicols}
\end{document}
